\documentclass[12pt,a4paper]{article}
\usepackage[utf8]{inputenc}
\usepackage{amsmath, amssymb}
\usepackage{geometry}
\usepackage{setspace}
\usepackage{graphicx}

\geometry{margin=2.5cm}
\setstretch{1.2}

\title{Taller de Matemáticas Básicas I}
\author{José Luciano Mejía Arias}
\date{}

\makeatletter
\renewcommand{\maketitle}{
  \begin{flushright}
    \includegraphics[width=3cm]{../../logo_unal.png}
  \end{flushright}
  \begin{center}
    {\LARGE \textbf{\@title}}\\[0.3cm]
    {\large \@author}
  \end{center}
}
\makeatother

\begin{document}
\maketitle

\section*{Instrucciones}
Resuelva cada uno de los siguientes ejercicios mostrando todo el proceso. Presente su trabajo de manera ordenada y clara.

\section{Jerarquía de Operaciones}
Resuelva las siguientes expresiones respetando la jerarquía de operaciones:

\begin{enumerate}
    \item $5 + 3 \times 2^2 - 8 \div 4$
    \item $\left(12 - 3\right) \times 2 + 5^2 \div 5$
    \item $\frac{15 + 3 \times 4}{9 - 2^2}$
    \item $2 \times \left[8 - 3 \times (4 - 1)\right] + 5$
    \item $\frac{6 \times 4 - 2^3}{3 + 1} + 7 \times 2$
\end{enumerate}

\section{Homogenización}
Simplifique las siguientes expresiones algebraicas combinando términos semejantes:

\begin{enumerate}
    \item $5x + 3y - 2x + 7y - x$
    \item $4a^2 - 3ab + 2b^2 + 5ab - a^2 + 3b^2$
    \item $\frac{2x}{3} + \frac{x}{4} - \frac{5x}{6}$
    \item $3(x + 2) - 2(x - 1) + 4(2x + 3)$
    \item $\frac{x + 1}{2} + \frac{2x - 3}{4} - \frac{x - 2}{3}$
\end{enumerate}

\section{Factorización}
Factorice completamente las siguientes expresiones:

\begin{enumerate}
    \item $x^2 + 5x + 6$
    \item $4x^2 - 9$
    \item $x^2 - 6x + 9$
    \item $2x^2 + 7x + 3$
    \item $x^3 - 8$
    \item $6x^2 - 13x - 5$
    \item $x^2 + 6x + 9 - y^2$
    \item $ax + ay + bx + by$
\end{enumerate}

\section{Ecuaciones}
Resuelva las siguientes ecuaciones:

\begin{enumerate}
    \item $3x + 7 = 22$
    \item $\frac{2x - 3}{4} = \frac{x + 1}{2}$
    \item $x^2 - 5x + 6 = 0$
    \item $2x^2 + 3x - 2 = 0$
    \item $\frac{1}{x - 1} + \frac{2}{x + 1} = 3$
    \item $\sqrt{2x + 5} = 3$
    \item $|3x - 4| = 8$
    \item $(x - 2)^2 = 16$
\end{enumerate}

\section{Problemas de Aplicación}
Resuelva los siguientes problemas planteando y resolviendo la ecuación correspondiente:

\begin{enumerate}
    \item La suma de dos números consecutivos es 47. ¿Cuáles son esos números?
    
    \item Un rectángulo tiene un perímetro de 36 metros. Si el largo es el doble del ancho, encuentre las dimensiones del rectángulo.
    
    \item Pedro tiene el triple de la edad de Juan. Si dentro de 8 años la suma de sus edades será 56 años, ¿cuál es la edad actual de cada uno?
    
    \item Un tanque de agua se llena con dos llaves. La primera llave puede llenar el tanque en 6 horas y la segunda en 4 horas. ¿En cuánto tiempo se llenará el tanque si se abren ambas llaves simultáneamente?
    
    \item Un comerciante compró cierto número de artículos por \$240. Si cada artículo hubiera costado \$2 menos, habría podido comprar 4 artículos más por el mismo dinero. ¿Cuántos artículos compró?
\end{enumerate}

\end{document}
